\documentclass[11pt,a4paper]{article}
\usepackage[margin=23mm]{geometry}
\usepackage[T1]{fontenc}
\usepackage{lmodern}
\usepackage{amsmath,amssymb,longtable,booktabs,array}
\usepackage[hidelinks]{hyperref}
\usepackage{listings}
\lstset{basicstyle=\ttfamily\small,breaklines=true,columns=fullflexible}
\setlength{\emergencystretch}{3em}
\title{MEIKURAL Version 2\\One recording, one Sa, unchanged spectrogram CNN}
\author{Implementation and mathematical reference}
\date{}
\begin{document}
\maketitle
\tableofcontents

\section{Purpose and controlled change}
Version 2 estimates Sa once from an entire supplied recording and reuses that
frequency for every 30-second window. Version 1 estimates Sa independently per
window. This is the intended experimental difference: the CNN architecture,
window geometry, per-window gain normalization, STFT, relative-pitch grid,
compression, loss, optimizer, regularization, and training defaults are unchanged.
The model still classifies Shree, Jog, Miyan Malhar, Lalit, Todi, Marwa, and Bhairavi.

All new implementation lives under \path{src/ver2}. Unchanged signal-processing,
CNN, configuration validation, and plotting primitives are imported from Version 1
rather than independently reimplemented. Version 1 code, checkpoints, and features
remain usable with Version 1 entry points. Separate Version 2 orchestration selects
the recording-level feature pipeline and rejects Version 1 artifacts.

\section{Inputs, artifacts, and future execution}
\path{src/ver2/config.json} is the single editable configuration. Its default
training source is \path{data/ver2/set}, feature cache is \path{data/ver2/features},
and model run is \path{models/ver2}. Relative paths resolve against the repository
root. Populate or expand these sources before starting; the implementation does
not copy the Version 1 dataset, create features, or train automatically. You may
point the dataset setting at an existing raw-audio collection for a controlled
comparison, while keeping feature and run output directories separate.

Place training performances beneath the seven configured class folders. Expanding
the number of performances does not change the network. Expanding the class set
would change the output task and is not part of this version. All eligible training
recordings are used, with no internal validation or test split. Optional annotations
are diagnostics only. Batch inference defaults to the existing shared
\path{data/test_ver1_handpicked} collection, with ground truth supplied by its raag
folders. The notebook's test root remains editable; reports use Version 2 paths.
Audio is the only inference input.

The optional dataset-selection command reuses Version 1's UUID-based planner and
atomic complete-folder copy implementation. It reads the selected classes and
destination from the Version 2 JSON, scans every Saraga recording's metadata,
and copies all recordings with exactly one matching raag label, including their
annotations. There is no artist or per-class sampling limit. Unlabelled and
multi-raag recordings are excluded. Identical destination folders are skipped;
conflicting contents are not overwritten. A dry run of the current local source
finds all 25 matching recordings (the same full selection as Version 1).

\subsection{Commands to run later}
Use the existing project virtual environment and unified requirements. CUDA
training remains explicit; automatic fallback is available by setting the training
device to \texttt{auto}, and explicit CPU mode remains supported.
\begin{lstlisting}
cd /home/mushu-cislab/GitHub/MEIKURAL
source .venv/bin/activate
python3 -m pip install -r requirements.txt
# Training: data/ver2/set; shared test audio: data/test_ver1_handpicked.
# These commands are user-launched, not run by the implementation:
python3 -m src.ver2.preprocessing.copy_raag_folders --dry-run
python3 -m src.ver2.preprocessing.copy_raag_folders
python3 -m src.ver2.preprocessing.prepare --diagnose-tonic
python3 -m src.ver2.training.train
python3 -m src.ver2.inference.predict \
    /path/to/audio.mp3 models/ver2/latest.pt --device auto
\end{lstlisting}
No trained Version 2 model is provided before that training step. Version 1
weights must not be renamed and used as Version 2 weights: they learned from a
different feature-generation procedure.

\section{Decoding and full-recording tonic estimation}
\subsection{Real audio and gain}
Decode all audio, average channels arithmetically, and resample to feature rate
$f_s=22050$ Hz with the unchanged decoder. Native decoding uses stateful soxr HQ;
unsupported codecs use FFmpeg with explicit arithmetic channel averaging.
Missing, empty, corrupt, or nonfinite audio is an error. A recording is decoded
once; files are processed sequentially.

Let the complete real recording be $x_0,\ldots,x_{N-1}$, with no artificial
padding. For both global tonic input and individual chunk normalization, use
\[
r(x)=\sqrt{\frac1N\sum_n x_n^2},\qquad p(x)=\max_n|x_n|,
\]
\[
g(x)=\min\left(\frac{10^{R/20}}{r(x)},10^{G/20},\frac{P}{p(x)}\right).
\]
Defaults are target $R=-20$ dBFS, amplification limit $G=20$ dB, and peak ceiling
$P=0.99$. Reject near-silence below RMS $10^{-5}$; no clipping or excessive silence
amplification is used. RMS accumulates in float64; normalized samples are float32.

\subsection{One estimator call for the entire performance}
Normalize the entire recording with this gain rule and resample it from the
feature rate to the estimator's 44100 Hz rate. Pass the complete resulting vector
to Essentia's \texttt{TonicIndianArtMusic} exactly once. The tail is included.
The input is not a selection of windows, a median of window estimates, or the
first 30 seconds. Full-recording RMS normalization prepares only the estimator
input; original decoded samples remain available for independent chunk gain
normalization.

The estimator uses harmonic multipitch salience, a pitch histogram accumulated
across frames, and its built-in decision tree to select a likely tonic
\cite{tonic,essentia}. The algorithm is reset between recordings because its frame
cutter retains state. The same 65--500 Hz search range and all other parameters
from Version 1 remain unchanged; see the generated tables.

Writing $\mathcal T$ for this unchanged algorithm and $\mathcal R$ for resampling:
\[
\widehat f_{\mathrm{Sa}}=\mathcal T\bigl(\mathcal R(g(x)x)\bigr),\qquad
\widehat f_{\mathrm{Sa},j}=\widehat f_{\mathrm{Sa}}\quad\text{for every chunk }j.
\]
Reject nonfinite or out-of-range estimates. If no tonic can be estimated, every
window of that recording is reported unusable; there is no per-window fallback,
annotation fallback, or guessed tonic. A valid estimator output is not a
calibrated confidence measurement and can still be wrong.

\subsection{Memory and scope of available information}
The standard Essentia interface accepts an entire audio vector, so the decoded
recording is held in RAM. Mono float32 at the feature rate uses about 5.3 MB per
minute, before normalization, resampling, the estimator's internal buffers, and
temporary copies. Peak memory is substantially larger. Very long performances
may require a streaming implementation in a future version. No undocumented
subsampling or truncated estimation is substituted to reduce this cost.

Version 2 is offline inference: later parts of a recording can inform the Sa used
for earlier windows. A standalone 30-second input estimates Sa from those supplied
30 seconds. It will not necessarily get the same estimate as an excerpt selected
from a full song. The interactive notebook explicitly uses the full loaded song.

\section{Unchanged window construction and STFT image}
\subsection{Windows and padding}
At rate $f_s$, let $L=30f_s$, $S=15f_s$, and $a_0=-15f_s$. Generate starts
$a_j=a_0+jS$ until the first window reaches or crosses $N$:
\[
J=\max\left(0,\left\lceil\frac{N-L-a_0}{S}\right\rceil\right),\qquad j=0,\ldots,J.
\]
Keep real samples on $[\max(0,a_j),\min(N,a_j+L))$, with left padding
$\max(0,-a_j)$ and right padding $\max(0,a_j+L-N)$.
Exact-boundary EOF does not create an extra window. Per-window gain is measured
from real samples only and applied before padding. Silent windows are skipped,
but their metadata retains the recording's known Sa.

\subsection{STFT definition and resolution}
For padded normalized window $y_j$, FFT size $F$, centered analysis window $w$,
and hop $H$:
\[
Z_j[k,t]=\sum_{n=0}^{F-1}\widetilde y_j[tH+n-F/2]\,w[n]e^{-2\pi i kn/F},
\qquad M_j[k,t]=|Z_j[k,t]|.
\]
Defaults are a periodic Hann window of length $F=8192$, hop $H=512$, and zero
padding for centered frames. Librosa's unnormalized forward STFT convention is
unchanged. Frequencies are $kf_s/F$ and frame centers $tH/f_s$.
\[
T=1+\left\lfloor\frac{L+2\lfloor F/2\rfloor-F}{H}\right\rfloor=1292.
\]
FFT spacing is about 2.692 Hz, window duration about 371.5 ms, and hop about
23.22 ms. The Hann window's frequency discrimination limits actual pitch detail;
interpolating to a fine pitch grid cannot create missing spectral resolution.

\subsection{Fixed Sa remapping for every chunk}
Use the same cents grid as Version 1:
\[
c_i=-1200+25i,\quad i=0,\ldots,239,\qquad
f_i=\widehat f_{\mathrm{Sa}}2^{c_i/1200}.
\]
The grid ends at 4775 cents, retains separate octaves, and always uses the same
recording-level $\widehat f_{\mathrm{Sa}}$. For $u_i=f_iF/f_s$,
$k_i=\lfloor u_i\rfloor$, and $\alpha_i=u_i-k_i$:
\[
A_j[i,t]=(1-\alpha_i)M_j[k_i,t]+\alpha_iM_j[k_i+1,t].
\]
The configured maximum tonic and grid must remain below Nyquist. Time positions
are unchanged. This is frequency-axis resampling, not changing audio pitch or
separating vocals. Harmonics, accompaniment, and recording characteristics remain.

For $m_j=\max A_j$, numerical floor $\epsilon$, and dynamic range $B=80$ dB,
return zeros if $m_j\le\epsilon$; otherwise:
\begin{align*}
D_j[i,t]&=20\log_{10}\frac{\max(A_j[i,t],\epsilon)}{m_j},\\
X_j[i,t]&=\frac{\operatorname{clip}(D_j[i,t],-B,0)+B}{B}.
\end{align*}
The matrix is $240\times1292$ in $[0,1]$. Store float16 NumPy features; convert to
float32 for CNN input. Inference repeats that quantization for consistency. No
plot axes, color mapping, image-file compression, or resizing enter the network.

\section{Identical CNN and training mathematics}
The CNN class is imported directly from Version 1. Input is $(B,1,240,1292)$.
Biased same-padded stride-one convolutions use kernels $(3,7)$, $(5,5)$, and
$(3,3)$ with output channels 16, 32, and 64. Each convolution is followed by ReLU
and floor-mode max pooling with kernel and stride $(2,4)$. A convolution has
$C_o(C_i k_h k_t+1)$ parameters; pooling and ReLU have none.

After the pools, dimensions are $(16,120,323)$, $(32,60,80)$, and $(64,30,20)$.
Adaptive average pooling to $(12,1)$ averages time and retains 12 pitch regions.
Adaptive output bin $j$ averages input indices
$\lfloor jI/O\rfloor$ through $\lceil(j+1)I/O\rceil-1$; bins can overlap.
Flatten to 768 values, apply training-only dropout 0.25 with retained entries
scaled by $1/0.75$, then a biased 768-to-7 linear classifier.

The convolution parameter counts are 352, 12832, and 18496; the classifier has
5383, for a total of 37063. For receptive fields, initialize $r_0=j_0=1$ and use
\[
r_\ell=r_{\ell-1}+(k_\ell-1)j_{\ell-1},\qquad
j_\ell=j_{\ell-1}s_\ell.
\]
After the third pool, receptive fields are $26\times118$ bins/frames with jumps
$8\times64$. The time support is about $(117H+F)/f_s=3.09$ seconds before
adaptive pooling. The CNN aggregates local patterns rather than explicitly
modeling the complete order of a 30-second phrase.

Let $n_c$ be usable windows for class $c$, $N_w=\sum_c n_c$, and $C=7$.
Use $w_c=N_w/(Cn_c)$ and weighted cross-entropy on logits $z$:
\[
p_{b,c}=\frac{e^{z_{b,c}}}{\sum_k e^{z_{b,k}}},\qquad
\mathcal L=-\frac{\sum_b w_{y_b}\log p_{b,y_b}}{\sum_b w_{y_b}}.
\]
AdamW maintains $m_t=\beta_1m_{t-1}+(1-\beta_1)g_t$ and
$v_t=\beta_2v_{t-1}+(1-\beta_2)g_t^2$, with bias-corrected moments and update
\[
\theta_t=(1-\eta\lambda)\theta_{t-1}
-\eta\frac{\widehat m_t}{\sqrt{\widehat v_t}+\epsilon_{\mathrm{Adam}}}.
\]
Defaults remain batch 8, 30 epochs, learning rate $10^{-3}$, decay $10^{-4}$,
betas $(0.9,0.999)$, epsilon $10^{-8}$, and seed 42. Use float32, shuffled
batches, no augmentation, no scheduler, and no early stopping. Record training
loss/accuracy without calling them validation metrics. No optimizer steps or
training epochs have been run for Version 2 during implementation.

\section{Generated settings and exact model dimensions}
Regenerate this reference from the Version 2 JSON and real model modules with
\texttt{python3 -m src.ver2.documentation.generate}.
% Generated from config.json; do not edit numerical defaults here.
\subsection{Configured settings}
\begin{longtable}{p{0.49\linewidth}p{0.44\linewidth}}
\textbf{Setting} & \textbf{Value} \\\hline\endhead
paths.dataset & "data/ver2/set" \\
paths.features & "data/ver2/features" \\
paths.run & "models/ver2" \\
classes & Shree, Jog, Miyan Malhar, Lalit, Todi, Marwa, Bhairavi \\
audio.sample\_rate & 22050 \\
audio.window\_seconds & 30 \\
audio.stride\_seconds & 15 \\
audio.first\_start\_seconds & -15 \\
audio.rms\_dbfs & -20 \\
audio.max\_gain\_db & 20 \\
audio.peak\_ceiling & 0.99 \\
audio.silence\_rms & 1e-05 \\
audio.decode\_block\_samples & 65536 \\
tonic.sampleRate & 44100 \\
tonic.minTonicFrequency & 65 \\
tonic.maxTonicFrequency & 500 \\
tonic.frameSize & 2048 \\
tonic.hopSize & 512 \\
tonic.binResolution & 10 \\
tonic.numberHarmonics & 20 \\
tonic.harmonicWeight & 0.85 \\
tonic.magnitudeCompression & 1 \\
tonic.magnitudeThreshold & 40 \\
tonic.numberSaliencePeaks & 5 \\
tonic.referenceFrequency & 55 \\
spectrogram.n\_fft & 8192 \\
spectrogram.win\_length & 8192 \\
spectrogram.hop\_length & 512 \\
spectrogram.window & "hann" \\
spectrogram.center & true \\
spectrogram.min\_cents & -1200 \\
spectrogram.max\_cents\_exclusive & 4800 \\
spectrogram.step\_cents & 25 \\
spectrogram.dynamic\_range\_db & 80 \\
spectrogram.epsilon & 1e-12 \\
model.channels & [16, 32, 64] \\
model.kernels & [[3, 7], [5, 5], [3, 3]] \\
model.pool & [2, 4] \\
model.adaptive\_pool & [12, 1] \\
model.dropout & 0.25 \\
training.device & "cuda" \\
training.epochs & 30 \\
training.batch\_size & 8 \\
training.learning\_rate & 0.001 \\
training.weight\_decay & 0.0001 \\
training.betas & [0.9, 0.999] \\
training.epsilon & 1e-08 \\
training.seed & 42 \\
training.workers & 0 \\
inference.batch\_size & 8 \\
preprocessing.version & 2 \\
preprocessing.tonic\_scope & "recording" \\
\end{longtable}
\subsection{Derived CNN dimensions}
Feature shape: (240, 1292). Trainable parameters: 37,063.
\begin{longtable}{llll}
\textbf{Layer} & \textbf{Output (C,H,W)} & \textbf{Parameters} & \textbf{RF (H,W)} \\\hline\endhead
Input & 1 x 240 x 1292 & 0 & 1 x 1 \\
Conv2d 1 & 16 x 240 x 1292 & 352 & 3 x 7 \\
MaxPool2d 1 & 16 x 120 x 323 & 0 & 4 x 10 \\
Conv2d 2 & 32 x 120 x 323 & 12,832 & 12 x 26 \\
MaxPool2d 2 & 32 x 60 x 80 & 0 & 14 x 38 \\
Conv2d 3 & 64 x 60 x 80 & 18,496 & 22 x 70 \\
MaxPool2d 3 & 64 x 30 x 20 & 0 & 26 x 118 \\
AdaptiveAvgPool2d & 64 x 12 x 1 & 0 & aggregated \\
Flatten / Dropout & 768 & 0 & aggregated \\
Linear & 7 & 5,383 & aggregated \\
\end{longtable}


\section{Caches, checkpoints, and failures}
The Version 2 feature fingerprint hashes format 2, preprocessing version 2,
recording-level tonic scope, and the audio, tonic, and spectrogram settings.
Version 1 caches and checkpoints are incompatible
even though tensor dimensions match. The manifest includes source identities,
sizes/mtimes, configuration, a single tonic per recording, and window-level
timestamps, padding, tonic scope, features, and failures. A near-silent chunk does
not cause Sa to be re-estimated. Changed sources or preprocessing settings require
a new feature-cache directory.

Save checkpoints atomically after each future training epoch with model weights,
optimizer state, epoch, ordered classes, full config, feature fingerprint, and
package versions. Save training history and config alongside the checkpoint.
Existing checkpoints are protected against overwrite; automatic resume is not
implemented. The Version 2 loader rejects a Version 1 config before prediction,
validates the fingerprint, and restores saved architecture/preprocessing rather
than trusting current editable defaults.

Default training requires CUDA, failing clearly when unavailable. Explicit CPU
and automatic fallback remain available. Inference defaults to automatic device
selection and can load CUDA-trained weights on CPU. GPU memory errors propagate
without silently restarting training. Source audio is not modified. Small selected
Version 2 checkpoint/config/history files are eligible for Git, just like Version 1.

\section{Inference notebooks and evaluation}
\begin{itemize}
\item \path{notebooks/05_ver2_inference.ipynb}: raw-audio inference, fixed recording
Sa, seven-score timeline, ranked overall scores, and failure metadata.
\item \path{notebooks/06_ver2_interactive_clip.ipynb}: load the full song, estimate
Sa once, listen to and classify chosen passages using that same Sa. Changing
either file path requires loading again; moving the slider never estimates tonic.
\item \path{notebooks/07_ver2_handpicked_batch.ipynb}: per-performance combined
timeline/ranking figures and final song-level and chunk-level confusion matrices.
\end{itemize}
For successful windows with real durations $d_j$, keep the same overall ranking:
\[
\overline p_c=\frac{\sum_j d_jp_{j,c}}{\sum_j d_j}.
\]
This reduces padded-edge influence, not overlap correlation. A selected interactive
clip instead has one prediction and no whole-recording score average; its Sa still
comes from the full recording. A supplied short file has only its own content
available for tonic estimation.

Batch ground truth is the first raag folder below the test root. Preserve the
existing spelling mappings, including \texttt{sree} to Shree, while keeping unknown
labels separate. In each raw-count confusion matrix,
$C_{ij}=\sum_n\mathbf 1[y_n=i\wedge\widehat y_n=j]$.
The song matrix counts one final argmax per recording. The chunk matrix counts
every successful window's argmax once, including padded edges, against the parent
folder label. Exclusions are reported; failed recordings receive error figures.
PNG reports are saved under \path{outputs/ver2/handpicked}. Matrix cells reuse
retained labels and do not repeat inference.

\section{Verification and limitations}
Synthetic-audio tests verify one complete-recording estimator call, identical Sa
across all windows, tails and padding, no per-window fallback after a global
failure, versioned cache/checkpoint isolation, identical model architecture and
hyperparameters, and fixed tonic during interactive selections. Notebook code is
validated without running full-dataset inference. Dummy checkpoint round trips do
not train the model. Verification results are recorded in
\path{docs/ver2_verification.md}.

Full-recording estimation eliminates changes caused by separately selecting Sa
for adjacent windows. It does not prove the selected Sa is correct. Octave and
non-octave mistakes can now misalign an entire recording consistently. The fixed
reference assumes tonic is stable; medleys or deliberate tonic changes violate
that assumption. Timbre, accompaniment, and artist shortcuts can still dominate
classification. More overlapping windows do not create independent performers.
Scores are not calibrated probabilities, and unknown-raag rejection is absent.
Full preparation, training, and external accuracy evaluation are deliberately
deferred until the expanded datasets are supplied.

\begin{thebibliography}{9}
\bibitem{tonic} J. Salamon, S. Gulati, and X. Serra.
\emph{A Multipitch Approach to Tonic Identification in Indian Classical Music}.
ISMIR, 2012.
\bibitem{essentia} Essentia. \emph{TonicIndianArtMusic algorithm reference}.
\url{https://essentia.upf.edu/reference/std_TonicIndianArtMusic.html}.
\end{thebibliography}
\end{document}
